%---
% title: "Gửi Tôi Của Năm Mười Sáu Tuổi: Giá Như Khi Ấy Đại Số Tuyến Tính Được Viết Thế Này. Kì 4 - Ma trận nghịch đảo"
% date: 2026-09-27
% tags: [dai-so-tuyen-tinh, thpt, ma-tran, nghich-dao]
% summary: "Kì 4 chuỗi đại số tuyến tính cho THPT: tại sao không chia được cho ma trận, ma trận nghịch đảo là gì và vì sao là duy nhất, định thức 2×2 như một diện tích, giải hệ bằng nghịch đảo, cùng bốn luật đại số hỏng với ma trận."
%---
\documentclass[12pt,a4paper]{article}

% Khối khai báo dưới đây CHỈ dùng khi biên dịch PDF trực tiếp
% (Overleaf, pdfLaTeX). Blog không cần nó: Pandoc tự bỏ qua phần này.
\usepackage[utf8]{inputenc}
\usepackage[T5]{fontenc}
\usepackage[vietnamese]{babel}
\usepackage{amsmath,amssymb}
\usepackage[a4paper,margin=2.2cm]{geometry}

\begin{document}

Kì trước dừng lại ở một câu hỏi. Với ma trận, ta đã có phép cộng, phép trừ và phép nhân, nhưng còn thiếu phép chia.

Với số thực thì không có gì phải bàn, vì mọi \(a\neq0\) đều có số nghịch đảo.

Với hai ma trận vuông cùng cỡ thì khó hơn. Kì trước ta đã thấy \(AB\) và \(BA\) đều tồn tại nhưng có thể khác nhau, nên câu hỏi ``chia \(A\) cho \(B\)'' không có nghĩa rõ ràng.

Vì vậy kì này ta hỏi theo hướng khác. Thay vì hỏi chia thế nào, ta tìm thứ có thể \textbf{hoàn tác} phép nhân.

\subsection{Ma trận đơn vị}

Muốn có một ma trận nhân vào mà không làm thay đổi gì, ta thử đặt số \(1\) trên đường chéo chính và số \(0\) ở mọi chỗ khác. Ma trận cỡ \(n\times n\) như vậy gọi là \textbf{ma trận đơn vị} (\emph{identity matrix}), kí hiệu \(I_n\).
$$
I_2=\left[ \begin{array}{cc} 1 & 0\\ 0 & 1 \end{array} \right],\qquad
I_3=\left[ \begin{array}{rrr} 1 & 0 & 0\\ 0 & 1 & 0\\ 0 & 0 & 1 \end{array} \right].
$$
Nhân nó vào bên trái hay bên phải một ma trận, ta đều nhận lại đúng ma trận ấy. Chẳng hạn với \(A=\left[ \begin{array}{rrr} 1 & 2 & 4\\ 2 & 6 & 0 \end{array} \right]\).
$$
I_2A=\left[ \begin{array}{rr} 1 & 0\\ 0 & 1 \end{array} \right]\left[ \begin{array}{rrr} 1 & 2 & 4\\ 2 & 6 & 0 \end{array} \right]=\left[ \begin{array}{rrr} 1 & 2 & 4\\ 2 & 6 & 0 \end{array} \right]=A,
$$
$$
AI_3=\left[ \begin{array}{rrr} 1 & 2 & 4\\ 2 & 6 & 0 \end{array} \right]\left[ \begin{array}{rrr} 1 & 0 & 0\\ 0 & 1 & 0\\ 0 & 0 & 1 \end{array} \right]=\left[ \begin{array}{rrr} 1 & 2 & 4\\ 2 & 6 & 0 \end{array} \right]=A.
$$
Lý do nằm ngay trong luật hàng--cột ở kì trước. Hàng thứ \(i\) của \(I_n\) chỉ có đúng một số khác không, là số \(1\) ở vị trí \(i\). Nhân hàng ấy với cột \(j\) của \(A\), tổng chỉ còn một số hạng, và số hạng đó là \(a_{ij}\).

Ma trận đơn vị cũng là duy nhất. Giả sử một ma trận \(J\) cũng để yên mọi ma trận khi nhân vào. Lấy \(A\) là ma trận chỉ có số \(1\) ở vị trí \((i,i)\), ta được \(j_{ii}=1\); lấy \(A\) là ma trận chỉ có số \(1\) ở vị trí \((i,j)\) với \(i\neq j\), ta được \(j_{ij}=0\). Vậy \(J=I_n\).

Vì thế với mọi ma trận vuông, ta có thể đặt \(A^0=I\), giống như với số thực.

\subsection{Nghịch đảo, và nó là duy nhất}

\textbf{Định nghĩa 1.} Nếu \(A\) là ma trận vuông và tồn tại ma trận \(B\) cùng cỡ sao cho
$$
AB=BA=I, \tag{1}
$$
thì \(A\) được gọi là \textbf{ma trận khả nghịch} (\emph{invertible}), và \(B\) là một \textbf{nghịch đảo} của \(A\).

Đổi vai hai ma trận cho nhau, điều kiện (1) vẫn thế, tức \(BA=AB=I\). Nên \(A\) cũng là nghịch đảo của \(B\), và ta gọi chúng là hai ma trận nghịch đảo của nhau.

Nếu \(A\) khả nghịch thì nghịch đảo của nó là duy nhất, và ta kí hiệu nó là \(A^{-1}\). Chứng minh khá ngắn.

\textbf{Định lý 1.} Nếu \(B\) và \(C\) đều là nghịch đảo của \(A\) thì \(B=C\).

\emph{Chứng minh.} Vì \(BA=I\), nhân hai vế với \(C\) ở bên phải.
$$
(BA)C=IC.
$$
Chuyển ngoặc theo luật kết hợp (sẽ nói ở mục sau), vế trái thành \(B(AC)=BI=B\), còn vế phải là \(C\). Vậy \(B=C\). \(\square\)

Vì nghịch đảo là duy nhất, ta không phải bận tâm đang dùng nghịch đảo nào, và viết gọn
$$
AA^{-1}=A^{-1}A=I. \tag{2}
$$

\begin{quote}
Cảnh báo. Kí hiệu \(A^{-1}\) không có nghĩa là \(\frac1A\). Chia cho một ma trận không phải là một phép tính được định nghĩa trong bài này. Kí hiệu ấy được đặt ra để chỉ đúng cái mà (2) mô tả, và không có gì hơn.
\end{quote}

\textbf{Ví dụ 1.} Với
$$
A=\left[ \begin{array}{rr} 2 & -5\\ -1 & 3 \end{array} \right],\qquad
B=\left[ \begin{array}{rr} 3 & 5\\ 1 & 2 \end{array} \right],
$$
ta tính
$$
AB=\left[ \begin{array}{rr} 2 & -5\\ -1 & 3 \end{array} \right]\left[ \begin{array}{rr} 3 & 5\\ 1 & 2 \end{array} \right]=\left[ \begin{array}{rr} 1 & 0\\ 0 & 1 \end{array} \right]=I,
$$
$$
BA=\left[ \begin{array}{rr} 3 & 5\\ 1 & 2 \end{array} \right]\left[ \begin{array}{rr} 2 & -5\\ -1 & 3 \end{array} \right]=\left[ \begin{array}{rr} 1 & 0\\ 0 & 1 \end{array} \right]=I.
$$
Vậy hai ma trận này là nghịch đảo của nhau.

\subsection{Ma trận nào không có nghịch đảo}

Không phải ma trận vuông nào cũng có một ma trận \(B\) thỏa (1). Xét
$$
A=\left[ \begin{array}{rrr} 1 & 4 & 0\\ 2 & 5 & 0\\ 3 & 6 & 0 \end{array} \right].
$$
Cột thứ ba của \(A\) toàn số \(0\). Gọi \(\mathbf{c}_1,\mathbf{c}_2,\mathbf{0}\) là ba cột của \(A\). Với mọi ma trận \(B\) cỡ \(3\times3\), công thức (3) ở kì trước cho
$$
BA=B[\,\mathbf{c}_1\ \mathbf{c}_2\ \mathbf{0}\,]=[\,B\mathbf{c}_1\ B\mathbf{c}_2\ \mathbf{0}\,].
$$
Cột thứ ba của \(BA\) là vector \(0\), nên \(BA\neq I\). Không có \(B\) nào thỏa (1).

Ma trận như thế gọi là \textbf{ma trận suy biến} (\emph{singular}).

Lập luận trên chỉ dùng công thức tính tích theo cột ở kì trước, không cần gì mới. Một ma trận vuông có một hàng hoặc một cột toàn số \(0\) thì luôn suy biến.

\subsection{Ma trận cỡ hai và định thức}

Xét
$$
A=\left[ \begin{array}{cc} a & b\\ c & d \end{array} \right].
$$
Thay vì nhận công thức có sẵn, ta tự tìm. Theo (1), cần một ma trận \(X\) sao cho \(AX=I\), tức là
$$
\left[ \begin{array}{cc} a & b\\ c & d \end{array} \right]\left[ \begin{array}{cc} x_{11} & x_{12}\\ x_{21} & x_{22} \end{array} \right]=\left[ \begin{array}{cc} 1 & 0\\ 0 & 1 \end{array} \right].
$$
Nhân ra, ta được bốn phương trình. Chúng tách thành hai hệ hai ẩn, mỗi hệ cho một cột của \(X\).

Hệ thứ nhất ứng với cột thứ nhất.
$$
\begin{cases} ax_{11}+bx_{21}=1\\ cx_{11}+dx_{21}=0. \end{cases}
$$
Nhân phương trình đầu với \(d\), phương trình sau với \(b\), rồi trừ.
$$
(ad-bc)x_{11}=d.
$$
Tương tự, nhân phương trình đầu với \(c\), phương trình sau với \(a\), rồi trừ theo chiều ngược lại.
$$
(ad-bc)x_{21}=-c.
$$

Với hệ thứ hai, ứng với cột thứ hai, làm y hệt thì được
$$
(ad-bc)x_{12}=-b,\qquad (ad-bc)x_{22}=a.
$$

Vậy nếu \(ad-bc\neq0\) thì
$$
X=\frac{1}{ad-bc}\left[ \begin{array}{rr} d & -b\\ -c & a \end{array} \right]. \tag{3}
$$
Bốn phương trình vừa giải nói rằng bốn phần tử của \(AX-I\) đều bằng \(0\), nên ma trận \(X\) ở trên thỏa \(AX=I\). Ta chưa kiểm \(XA=I\), nhưng Định lý 1 nói nghịch đảo là duy nhất, nên khi \(AX=I\) thì \(X\) là nghịch đảo của \(A\).

Ngược lại, nếu \(ad-bc=0\) thì bốn phương trình trên không giải được. Mục sau sẽ cho thấy khi đó ma trận ép cả mặt phẳng xuống một đường thẳng.

Biểu thức \(ad-bc\) xuất hiện ở cả bốn phương trình, nên ta đặt tên cho nó.

\textbf{Định nghĩa 2.} Với \(A=\left[ \begin{array}{cc} a & b\\ c & d \end{array} \right]\), ta gọi
$$
\det(A)=ad-bc \tag{4}
$$
là \textbf{định thức} của \(A\). Với ma trận vuông cỡ \(n\times n\), định thức sẽ được định nghĩa tổng quát sau này; với cỡ hai thì (4) là tất cả.

Cách tìm ở trên cho ngay kết luận. Ma trận \(A\) cỡ \(2\times2\) khả nghịch khi và chỉ khi \(\det(A)\neq0\), và khi đó (3) là công thức nghịch đảo.

\textbf{Ví dụ 2.} (a) Với \(A=\left[ \begin{array}{rr} 6 & 1\\ 5 & 2 \end{array} \right]\) thì \(\det(A)=6\cdot2-1\cdot5=7\neq0\), nên
$$
A^{-1}=\frac17\left[ \begin{array}{rr} 2 & -1\\ -5 & 6 \end{array} \right]
=\left[ \begin{array}{rr} \frac27 & -\frac17\\ -\frac57 & \frac67 \end{array} \right].
$$
(b) Với \(A=\left[ \begin{array}{rr} -1 & 2\\ 3 & -6 \end{array} \right]\) thì \(\det(A)=(-1)(-6)-2\cdot3=0\), nên \(A\) suy biến.

\subsection{Định thức đo diện tích}

Cách tìm ở trên cho ta công thức, nhưng chưa cho biết \(ad-bc\) có nghĩa gì. Kì trước còn bỏ ngỏ một câu hỏi tự nhiên hơn nhiều. Phép biến đổi mà ma trận \(A\) mô tả làm diện tích thay đổi bao nhiêu lần?

Với \(A=\left[ \begin{array}{cc} a & b\\ c & d \end{array} \right]\), hai cột của \(A\) là \((a,c)\) và \((b,d)\). Theo cách xếp cột ở kì trước, ma trận này biến vector \((x,y)\) thành \((ax+by,\;cx+dy)\), tức là thành tổng
$$
x(a,c)+y(b,d).
$$
Vì vậy ô vuông đơn vị ở góc phần tư thứ nhất bị kéo thành hình bình hành có hai cạnh \((a,c)\) và \((b,d)\).

Hình bình hành ấy có diện tích \(|ad-bc|\). Chọn \((a,c)\) làm đáy thì chiều cao là phần hướng lên trục \(Oy\) của cạnh \((b,d)\), và tích đáy với cao cho ra đúng \(|ad-bc|\).

\blogfig{\includegraphics{/images/nghich-dao-ma-tran/dien-tich-dinh-thuc.svg}}

Với \(A=\left[ \begin{array}{rr} 2 & 1\\ 0 & 1 \end{array} \right]\) thì \(\det(A)=2\). Hình bình hành có cạnh \((2,0)\) và \((1,1)\), diện tích gấp \(2\) lần ô vuông đơn vị.

Với \(A=\left[ \begin{array}{rr} 1 & 1\\ 1 & 1 \end{array} \right]\) thì \(\det(A)=0\), vì hai cột bằng nhau. Hai cạnh cùng hướng nên hình bình hành bẹp thành một đoạn thẳng, diện tích bằng \(0\).

Như vậy \(\det(A)=0\) có một nghĩa rất cụ thể, chứ không chỉ là một bước rút gọn. Ma trận ấy không giữ được mặt phẳng mà nén nó xuống một đường thẳng. Mọi điểm khác nhau trên mặt phẳng đều bị đưa về cùng một chỗ, và thông tin về vị trí bị mất.

Nó cũng giải thích vì sao \(\det(A)\neq0\) đi cùng với việc \(A\) khả nghịch. Ma trận nào không ép dẹt mặt phẳng thì còn đủ thông tin để dựng lại, nên còn tìm được ma trận hoàn tác nó.

\subsection{Giải hệ phương trình bằng nghịch đảo}

Với hệ hai ẩn
$$
\begin{cases} ax+by=u\\ cx+dy=v, \end{cases}
$$
viết lại thành \(A\mathbf{x}=\mathbf{b}\) với \(\mathbf{x}=\left( \begin{array}{c} x\\ y \end{array} \right)\) và \(\mathbf{b}=\left( \begin{array}{c} u\\ v \end{array} \right)\). Nhân hai vế với \(A^{-1}\) ở bên trái.
$$
A^{-1}A\mathbf{x}=A^{-1}\mathbf{b}
$$
nên \(I\mathbf{x}=A^{-1}\mathbf{b}\), tức là \(\mathbf{x}=A^{-1}\mathbf{b}\). Thay (3) vào.
$$
\left( \begin{array}{c} x\\ y \end{array} \right)
=\frac{1}{ad-bc}\left[ \begin{array}{rr} d & -b\\ -c & a \end{array} \right]\left( \begin{array}{c} u\\ v \end{array} \right),
$$
tức là
$$
x=\frac{du-bv}{ad-bc},\qquad y=\frac{av-cu}{ad-bc}. \tag{5}
$$

Đã có khử Gauss--Jordan từ hai kì trước, ta có cần thêm cách này không?

Với một hệ hai ẩn có hệ số là những con số cụ thể thì không. Khử nhanh hơn, và không phải chia gì cả. Nghịch đảo có lợi khi hệ số là chữ, như ở (5). Chạy khử trên chữ thì bước nào cũng phải mang theo bốn chữ \(a,b,c,d\) cùng một mẫu số, rất dễ nhầm. Còn (5) thì đã có sẵn.

Lợi thế lớn hơn nằm ở chỗ khác. Với hệ vuông \(n\) ẩn, giải hệ không còn là chạy một thuật toán mà chỉ là một phép nhân. Nghiệm là \(A^{-1}\mathbf{b}\). Cùng một ma trận \(A^{-1}\) biến mọi vế phải thành nghiệm. Những kì sau sẽ dùng điều đó, khi phải xử lý hàng trăm phương trình cùng lúc.

\subsection{Nghịch đảo của một tích}

Giả sử \(A\) và \(B\) đều khả nghịch. Khi đó tích \(AB\) cũng khả nghịch, và nghịch đảo của nó là \(B^{-1}A^{-1}\).

\textbf{Định lý 2.} Nếu \(A\) và \(B\) là hai ma trận vuông khả nghịch cùng cỡ thì \(AB\) khả nghịch và
$$
(AB)^{-1}=B^{-1}A^{-1}. \tag{6}
$$

\emph{Chứng minh.} Chỉ cần nhân thử theo cả hai chiều.
$$
(AB)(B^{-1}A^{-1})=A(BB^{-1})A^{-1}=AIA^{-1}=AA^{-1}=I.
$$
Chiều kia, $(B^{-1}A^{-1})(AB)=B^{-1}(A^{-1}A)B=B^{-1}IB=B^{-1}B=I$. Vậy (6) đúng. \(\square\)

Thứ tự đã bị đảo, y như công thức ở mục chuyển vị của kì trước. Hai ma trận khả nghịch nói chung không giao hoán, nên không thể đảo thứ tự trong (6).

\textbf{Ví dụ 3.} Với
$$
A=\left[ \begin{array}{rr} 1 & 2\\ 1 & 3 \end{array} \right],\qquad
B=\left[ \begin{array}{rr} 3 & 2\\ 2 & 2 \end{array} \right],
$$
ta có
$$
AB=\left[ \begin{array}{rr} 7 & 6\\ 9 & 8 \end{array} \right],\qquad
(AB)^{-1}=\frac12\left[ \begin{array}{rr} 8 & -6\\ -9 & 7 \end{array} \right]=\left[ \begin{array}{rr} 4 & -3\\ -\frac92 & \frac72 \end{array} \right].
$$
Mặt khác, \(\det(A)=1\) và \(\det(B)=2\), nên
$$
A^{-1}=\left[ \begin{array}{rr} 3 & -2\\ -1 & 1 \end{array} \right],\qquad
B^{-1}=\left[ \begin{array}{rr} 1 & -1\\ -1 & \frac32 \end{array} \right].
$$
Nhân theo đúng thứ tự (6).
$$
B^{-1}A^{-1}
=\left[ \begin{array}{rr} 1 & -1\\ -1 & \frac32 \end{array} \right]\left[ \begin{array}{rr} 3 & -2\\ -1 & 1 \end{array} \right]
=\left[ \begin{array}{rr} 4 & -3\\ -\frac92 & \frac72 \end{array} \right],
$$
đúng bằng \((AB)^{-1}\).

Nhân theo thứ tự ngược lại thì ra một ma trận khác.
$$
A^{-1}B^{-1}
=\left[ \begin{array}{rr} 3 & -2\\ -1 & 1 \end{array} \right]\left[ \begin{array}{rr} 1 & -1\\ -1 & \frac32 \end{array} \right]
=\left[ \begin{array}{rr} 5 & -6\\ -2 & \frac52 \end{array} \right]\neq (AB)^{-1}.
$$

\textbf{Định lý 3.} Nếu \(AB\) là ma trận suy biến thì \(A\) hoặc \(B\) là ma trận suy biến.

\emph{Chứng minh.} Phản chứng. Nếu cả \(A\) và \(B\) đều khả nghịch thì theo Định lý 2, \(AB\) khả nghịch, trái với giả thiết \(AB\) suy biến. \(\square\)

Nói cách khác, nhân hai ma trận khả nghịch không bao giờ ra một ma trận suy biến. Tích đã suy biến thì ít nhất một thừa số phải suy biến.

\subsection{Những luật còn giữ, và bốn luật không}

Ở kì trước, phép cộng hai ma trận và phép nhân ma trận với một số đều làm trên từng phần tử, còn mỗi phần tử của tích hai ma trận là một tổng các tích. Nghĩa là luật nào đúng với số thực mà chỉ dùng hai phép cộng và nhân thì vẫn đúng với ma trận.

Sách giáo trình liệt kê hơn hai mốt luật như vậy. Tất cả đều đúng theo cùng một cơ chế, nên xem cơ chế ấy một lần là đủ. Với ba ma trận \(A,B,C\) cùng cỡ và các số \(a,b\), ta có
$$
\begin{aligned}
A(B+C)_{ij} &= \sum_{k}a_{ik}(b_{kj}+c_{kj})\\
            &= \sum_{k}\bigl(a_{ik}b_{kj}+a_{ik}c_{kj}\bigr)\\
            &= (AB)_{ij}+(AC)_{ij}
\end{aligned}
$$
nên \(A(B+C)=AB+AC\). Ta chỉ phá ngoặc rồi gom lại, giống hệt với số thực. Các luật sau được kiểm tra theo đúng cách ấy.
$$
\begin{gathered}
A+B=B+A,\qquad A+(B+C)=(A+B)+C,\\
A(BC)=(AB)C,\qquad (B+C)A=BA+CA,\\
a(B+C)=aB+aC,\qquad (a+b)C=aC+bC,\\
a(bC)=(ab)C,\qquad a(BC)=(aB)C=B(aC).
\end{gathered}
$$

Có bốn luật không còn đúng, mỗi luật hỏng theo một kiểu.

\textbf{Phép nhân không giao hoán.} Kì trước đã chỉ ra điều này bằng hai ma trận cụ thể.

\textbf{Quy tắc giải bỏ hỏng.} Với số thực, \(ab=ac\) và \(a\neq0\) kéo theo \(b=c\). Với ma trận thì không.
$$
A=\left[ \begin{array}{rr} 0 & 1\\ 0 & 2 \end{array} \right],\qquad
B=\left[ \begin{array}{rr} 1 & 0\\ 3 & 4 \end{array} \right],\qquad
C=\left[ \begin{array}{rr} 2 & 0\\ 3 & 4 \end{array} \right].
$$
Ta có
$$
AB=\left[ \begin{array}{rr} 3 & 4\\ 6 & 8 \end{array} \right],\qquad
AC=\left[ \begin{array}{rr} 3 & 4\\ 6 & 8 \end{array} \right],
$$
trong khi \(B\neq C\). Nếu giải bỏ \(A\) ở hai vế, ta sẽ đi tới kết luận sai.

\begin{quote}
Chỗ hỏng nằm ở chỗ khác. Với số thực, điều kiện \(a\neq0\) là đủ để giải bỏ, vì \(a\) là một số. Nhân hai vế với cùng một số khác không thì không thể làm hai vế khác nhau trở thành bằng nhau. Với ma trận thì \(A\neq0\) không đủ. Ma trận \(A\) ở trên khác không, nhưng nó biến hai ma trận khác nhau thành cùng một ma trận.

Nên muốn một luật giải bỏ đúng với ma trận, ta phải đòi một điều mạnh hơn nhiều so với ``khác không''. Điều đó là khả nghịch.
\end{quote}

\textbf{Tích bằng không mà hai thừa số khác không.} Với số thực, \(ab=0\) thì \(a=0\) hoặc \(b=0\). Với ma trận, hai thừa số dưới đây đều khác không mà tích vẫn bằng không.
$$
\left[ \begin{array}{rr} 0 & 1\\ 0 & 2 \end{array} \right]\left[ \begin{array}{rr} 3 & 7\\ 0 & 0 \end{array} \right]=\left[ \begin{array}{rr} 0 & 0\\ 0 & 0 \end{array} \right].
$$
Hai cột của ma trận bên phải đều là bội của \((1,0)\), nên tích ra một ma trận mà mỗi hàng là một bội của hàng kia. Định lý 3 cho biết chuyện này không tránh được. Tích suy biến thì phải có thừa số suy biến, và ở đây ma trận bên phải suy biến vì có hàng thứ hai toàn số \(0\).

\textbf{Bình phương của tổng.} Với số thực, \((a+b)^2=a^2+2ab+b^2\), vì \(ab=ba\). Với ma trận thì ta chỉ viết được
$$
(A+B)^2=A^2+AB+BA+B^2,
$$
và chỉ rút gọn được thành \(A^2+2AB+B^2\) khi \(AB=BA\).

Với hai ma trận ở Ví dụ 3.
$$
A^2=\left[ \begin{array}{rr} 3 & 8\\ 4 & 11 \end{array} \right],\qquad
B^2=\left[ \begin{array}{rr} 13 & 10\\ 10 & 8 \end{array} \right],\qquad
AB=\left[ \begin{array}{rr} 7 & 6\\ 9 & 8 \end{array} \right],\qquad
BA=\left[ \begin{array}{rr} 5 & 12\\ 4 & 10 \end{array} \right].
$$
Nên
$$
(A+B)^2=\left[ \begin{array}{rr} 28 & 36\\ 27 & 37 \end{array} \right],\qquad
A^2+2AB+B^2=\left[ \begin{array}{rr} 30 & 30\\ 32 & 35 \end{array} \right].
$$
Hai ma trận này không bằng nhau.

\subsection{Lũy thừa và đa thức ma trận}

Với ma trận vuông, ta đặt các lũy thừa không âm theo cách quen thuộc
$$
A^0=I,\qquad A^1=A,\qquad A^2=AA,\qquad A^n=\underbrace{AA\cdots A}_{n\text{ ma trận}}.
$$
Và nếu \(A\) khả nghịch thì các lũy thừa âm là
$$
A^{-1}=A^{-1},\qquad A^{-2}=A^{-1}A^{-1},\qquad A^{-n}=\left(A^{-1}\right)^n.
$$
Các luật lũy thừa quen thuộc vẫn đúng, và lý do lần này đơn giản hơn thường lệ.
$$
A^{r}A^{s}=A^{r+s},\qquad (A^{r})^{s}=A^{rs}. \tag{7}
$$
Mọi lũy thừa của một ma trận đều là tích của chính nó với nhau, theo cùng một thứ tự. Chẳng hạn \(A^rA^s=AA\cdots A\,AA\cdots A\), gộp lại thành một dãy gồm \(r+s\) thừa số \(A\).

Với mọi đa thức \(p,q\), ta có \(p(A)q(A)=q(A)p(A)\). Muốn hiểu vì sao, trước hết cần định nghĩa đa thức ma trận.

Với \(A\) vuông và
$$
p(x)=a_0+a_1x+a_2x^2+\cdots+a_mx^m,
$$
ta đặt
$$
p(A)=a_0I+a_1A+a_2A^2+\cdots+a_mA^m. \tag{8}
$$
Hằng số \(a_0\) được thay bằng ma trận \(a_0I\), vì phép cộng một số với một ma trận chưa được định nghĩa. Cách viết này cũng khớp với quy ước \(A^0=I\) cho số hạng bậc không.

Các lũy thừa của \(A\) giao hoán với nhau, nên các đa thức của \(A\) cũng giao hoán với nhau. Giữa các đa thức của cùng một ma trận, ta có lại những điều mà hai ma trận tùy ý không có.

\textbf{Ví dụ 4.} Với \(p(x)=x^2-2x-3\) và
$$
A=\left[ \begin{array}{rr} -1 & 2\\ 0 & 3 \end{array} \right],
$$
ta có
$$
A^2=\left[ \begin{array}{rr} 1 & 4\\ 0 & 9 \end{array} \right],\qquad
A^2-2A-3I
=\left[ \begin{array}{rr} 1 & 4\\ 0 & 9 \end{array} \right]-\left[ \begin{array}{rr} -2 & 4\\ 0 & 6 \end{array} \right]-\left[ \begin{array}{rr} 3 & 0\\ 0 & 3 \end{array} \right]
=\left[ \begin{array}{rr} 0 & 0\\ 0 & 0 \end{array} \right].
$$
Vậy \(p(A)\) bằng ma trận không, dù \(p(x)=x^2-2x-3=(x-3)(x+1)\) không bằng không ở bất kì số nào. Điều đó không có nghĩa \(p\) bị sai. Ở đây ta trừ các ma trận, mà phép trừ làm trên từng phần tử; các phần tử của \(A^2\), của \(2A\) và của \(3I\) đều nằm đúng chỗ của chúng.

\subsection{Chuyển vị và nghịch đảo}

Ở mục chuyển vị của kì trước, ta đã có \((AB)^T=B^TA^T\). Ghép với (2) thì ra một đẳng thức nữa.

\textbf{Định lý 4.} Nếu \(A\) khả nghịch thì \(A^{T}\) cũng khả nghịch, và
$$
\left(A^{-1}\right)^{T}=\left(A^{T}\right)^{-1}. \tag{9}
$$

\emph{Chứng minh.} Chuyển vị hai vế của \(A^{-1}A=I\), rồi dùng \((AB)^T=B^TA^T\).
$$
A^{T}\left(A^{-1}\right)^{T}=\left(A^{-1}A\right)^{T}=I^{T}=I.
$$
Tương tự, $\left(A^{-1}\right)^{T}A^{T}=\left(AA^{-1}\right)^{T}=I^{T}=I$. Vậy $\left(A^{T}\right)^{-1}=\left(A^{-1}\right)^{T}$. \(\square\)

Với ma trận cỡ hai, điều này kiểm tra được bằng mắt. Ta có
$$
\left(\left[ \begin{array}{rr} a & b\\ c & d \end{array} \right]\right)^{T}
=\left[ \begin{array}{rr} a & c\\ b & d \end{array} \right],
$$
hai ma trận có cùng định thức, còn nghịch đảo tính theo (3) cũng chỉ đổi chỗ hai số \(-b\) và \(-c\).

\subsection{Ma trận lớn có khả nghịch không}

Ta đã có định nghĩa nghịch đảo, định lý tính duy nhất, và một công thức tường minh cho ma trận cỡ hai.

Nhưng công thức ấy chỉ dùng được cho cỡ hai. Một ma trận \(10\times10\) thì sao? Ở cỡ hai, ta chỉ việc sắp xếp lại bốn con số quanh đường chéo chính là quyết định được tất cả. Với một ma trận lớn, cần một cách khác để biết nó có khả nghịch hay không.

Cách ấy thật ra đã có sẵn trong tay mà ta chưa dùng tới. Ở kì trước, ta thấy hàng thứ \(i\) của \(A\mathbf{x}\) chính là phương trình thứ \(i\) của hệ. Vậy khử một hệ bằng phép biến đổi hàng, theo đúng nghĩa của nó, là biến đổi chính ma trận hệ số \(A\).

Khi khử \(A\) về dạng bậc thang rút gọn để đọc nghiệm, có phải ta đang khử \(A\) về ma trận đơn vị không? Nếu đúng vậy thì chỉ cần một phép biến đổi hàng duy nhất là biến \(A\) thành \(I\), và phép biến đổi ấy chính là nghịch đảo của \(A\).

Chứng minh điều đó cần thêm một việc, là gom cả một dãy phép biến đổi hàng thành một phép nhân duy nhất đặt ở bên trái. Việc ấy để dành cho kì sau.
\end{document}
